\scanpage{130}
\deleted{
Soit un elt $\unsure{f|\varphi}\in\operatorname{Ext}^i(C;G,S\unsure{\tau|\sigma})$, i.e. un hom de degré $i$ $G\mathrel{-}S\unsure{\tau|\sigma}$, d'où un hom de degré $i$ $Rf_*G\mathrel{-}Rf_*S\unsure{\tau|\sigma}$, donc il reste à définir un hom de degré $0$
\[
Rf_*S\unsure{\tau|\sigma}=Rf_*(\unsure{\tau|\sigma}\otimes\mu_n(1))\longrightarrow\unsure{\tau|\sigma}
\]
\[
R^{\unsure{i}}
\]
i.e. un élément de
\[
\operatorname{Ext}^2(k,Rf_*(\unsure{\tau|\sigma}\otimes\mu_n),\unsure{\tau})\mathrel{-}\quad---
\]
}
On a un hom \unsure{évident}
\[
\underline E_f(G,S\unsure{\tau|\sigma})\mathrel{-}E_f(Rf_*G,Rf_*S\unsure{\tau|\sigma})
\]
et si nous définissons un élément
\[
\gamma\in\deleted{\unsure{\operatorname{Ext}^0}}\inserted{\operatorname{Hom}^0}(Rf_*S\unsure{\tau|\sigma},\unsure{\tau|\sigma})
\]
nous en déduisons un \deleted{élément} hom
\[
E_f(Rf_*G,Rf_*S\unsure{\tau|\sigma})\mathrel{-}E_f(Rf_*G,\unsure{\tau|\sigma})
\]
et en composant, on trouve l'hom voulu
\[
E_f(G,S\unsure{\tau|\sigma})\longrightarrow E_f(Rf_*G,\unsure{\tau|\sigma})
\]
Ceci dit, le th. de dualité s'énonce

\uline{Th.} (i) Si $G$ est \inserted{un \deleted{\unsure{proalgébrique}} faisceau de Weil} fini [plus gén les $\unsure{H^i(G)}$ sont \deleted{\unsure{tous finis}}], \deleted{alors} plus gén est affine, et sans composantes abéliennes, alors
\[
E_f(G,S\unsure{\tau|\sigma})\xrightarrow{\sim}E(Rf_*G,\unsure{\tau|\sigma})
\]
\margin{\unsure{préciser}}
\figure{\begin{tikzpicture}\draw(0,0)--(0,1.2);\end{tikzpicture}}
(ii) Si $G$ est \inserted{de type fini} sans composantes abéliennes, alors $E_f(G,S\unsure{\tau|\sigma})\longrightarrow E(Rf_*G,\unsure{\tau|\sigma})$ est un isom mod $\theta$ [i.e. \unsure{les} $\underline E$ dans $\mathbb Q/\mathbb Z$ \deleted{\unsure{est un}} \unsure{se transforme} en un isom\ldots].
\mspage{129}

